\subsection{Implicit Function Theorem}

\begin{definition} \label{def:implicit-derivative-blocks} \hlt{Block Splitting of a Linear Map}\\
For $A\in\LL(\R^{n+m},\R^n)$ and $(\h,\mathbf{k})\in\R^n\times\R^m$, define $A_x\h\defeq A(\h,\ze)$ and $A_y\mathbf{k}\defeq A(\ze,\mathbf{k})$.\\
Then $A_x\in\LL(\R^n)$, $A_y\in\LL(\R^m,\R^n)$, and $A(\h,\mathbf{k})=A_x\h+A_y\mathbf{k}$.
\end{definition}

\begin{theorem} \label{thm:linear-implicit-function-theorem} \hlt{Linear Implicit Function Theorem}\\
If $A\in\LL(\R^{n+m},\R^n)$, $A_x$ invertible, then $\forall \mathbf{k}\in\R^m$ $\exists$ unique $\h\in\R^n$ s.t. $A(\h,\mathbf{k})=0$, that $\h=-(A_x)^{-1}A_y\mathbf{k}$.
\begin{proof}
Fix $\mathbf{k}\in\R^m$ and set $\h_0\defeq-(A_x)^{-1}A_y\mathbf{k}$.\\
By Definition \ref{def:implicit-derivative-blocks}, $A(\h_0,\mathbf{k}) =A_x\h_0+A_y\mathbf{k} =-A_y\mathbf{k}+A_y\mathbf{k} =0$, so a solution exists.\\
If $\h\in\R^n$ also satisfies $A(\h,\mathbf{k})=0$, then $A_x\h+A_y\mathbf{k}=0$, hence $A_x\h=-A_y\mathbf{k}$.\\
Applying $(A_x)^{-1}$ gives $\h=-(A_x)^{-1}A_y\mathbf{k}=\h_0$. Thus the solution exists and is unique.\\
\hlt{Structural Commentary}: Splitting $A$ separates the unknown variable from the parameter.\\
\hlt{Key Insight}: Invertibility of $A_x$ turns constraint $A(\h,\mathbf{k})=0$ into a uniquely solvable linear equation for $\h$.
\end{proof}
\end{theorem}

The implicit theorem below is the nonlinear counterpart of the block equation theorem. The invertible block is the set of variables one can solve for, just as in Theorem \ref{thm:linear-implicit-function-theorem}.

\begin{figure}[H]
  \centering
  \begin{tikzpicture}[
  scale=0.95,
  >=Stealth,
  every node/.style={font=\scriptsize},
  axis/.style={->, black!45, line width=0.5pt},
  window/.style={draw=black!55, fill=blue!4, line width=0.6pt},
  locus/.style={blue!70!black, line width=0.9pt},
  guide/.style={black!45, dashed, dash pattern=on 2.5pt off 2.5pt, line width=0.45pt},
  point/.style={circle, fill=black, inner sep=1.2pt}
]
  \draw[axis] (-0.15,0) -- (5.65,0) node[right] {$\mathbf{y}$};
  \draw[axis] (0,-0.15) -- (0,3.55) node[above] {$\mathbf{x}$};

  \draw[window] (2.05,0.92) rectangle (4.80,3.05);

  \coordinate (p) at (3.55,2.05);
  \draw[locus]
    (1.00,0.55) .. controls (1.70,1.10) and (2.35,1.42) .. (p)
    .. controls (4.02,2.38) and (4.70,3.08) .. (5.42,3.18);
  \node[blue!70!black, fill=white, inner sep=1pt] at (4.82,2.82) {$f(x,y)=0$};

  \draw[guide] (p) -- (3.55,0) node[below=2pt] {$b$};
  \draw[guide] (p) -- (0,2.05) node[left=2pt] {$g(b)$};
  \node[point] at (p) {};

  \node[above left=3pt] at (p) {$(g(b),b)$};
  \node[rotate=28, fill=white, inner sep=1pt] at (2.98,1.80) {$x=g(y)$};
\end{tikzpicture}

  \caption{Near a regular point, the zero locus is the graph of an implicit function.}
  \label{fig:implicit-locus-as-graph}
\end{figure}

\begin{theorem} \label{thm:implicit-function-theorem-rn-rm} \hlt{Implicit Function Theorem}\\
Let $\f:E\subseteq\R^{n+m}\to\R^n$ be $C^1$, $\f(\a,\b)=0$, $A=D\f(\a,\b)$, and let $A_x,A_y$ be the block maps of Definition \ref{def:implicit-derivative-blocks}. If $A_x$ is invertible, then near $\b$, the equation $\f(\x,\y)=0$ has unique local solution $\x=\g(\y)$, where $\g$ is $C^1$, $\g(\b)=\a$, $D\g(\b)=-(A_x)^{-1}A_y$.
\begin{proof}
Set $\Phi(\x,\y)=(\f(\x,\y),\y)$. Then $D\Phi(\a,\b)(\h,\mathbf{k})=(A_x\h+A_y\mathbf{k},\mathbf{k})$, whose kernel is zero since $A_x$ is invertible; hence $D\Phi(\a,\b)$ is invertible by Theorem \ref{thm:finite-dimensional-operator-invertible-equivalences}. Theorem \ref{thm:inverse-function-theorem-rn} makes $\Phi$ locally invertible.\\
Points with $\f(\x,\y)=0$ are exactly those sent to $(0,\y)$, so $\Phi^{-1}(0,\y)=(\g(\y),\y)$ gives the unique local solution. Differentiating $\f(\g(\y),\y)=0$ at $\b$ gives $A_xD\g(\b)+A_y=0$.\\
\hlt{Structural Commentary}: Add $\y$ as a coordinate, then invert.\\
\hlt{Key Insight}: Invertibility of the $x$-block is precisely the linear ability to solve for $x$.
\end{proof}
\end{theorem}

\begin{perspective} \label{per:rank-theorem-linear-model} \hlt{Rank Theorem Linear Model}\\
Theorem \ref{thm:quotient-by-kernel-is-range} reduces a linear map to its kernel--range split.\\
The constant-rank theorem is the local $C^1$ analogue after coordinate changes.
\end{perspective}
